% Resident narrative, city-centre case: Kingston upon Hull 029E (E01033104).
% Values: outputs/feature_local_coefficients.csv, outputs/geo_area_data_eval/n_4000/geo_coefficients.csv
% (area-data Geo design), outputs/geo_area_data_eval/spatial_patterns/lisa_lsoa.csv (LISA). Street names: OpenStreetMap (Overpass) within the official ONS LSOA 2021 boundaries.
Across Hull, the model's predictions depend mostly on each area's recent
crime history, especially the average over the last six months. This reliance
is strongest in the city centre and Old Town, and along the inner-city belt
through Dairycoates, Stepney and Sculcoates. Unemployment and long-term
sickness shares play a much smaller part, and features such as area size, road
density and junction density matter least.

Kingston upon Hull 029E (\texttt{E01033104}) is Hull's city centre and Old
Town: Queen Victoria Square, Whitefriargate, Jameson Street, Paragon Square and
Ferensway, Trinity Square, Market Place and Queens Dock Avenue, down to Humber
Street and Castle Street. Here the model leans on the six- and three-month crime
averages far more than anywhere else in Hull, while junction numbers,
long-term sickness and transit-stop counts matter least. The model is
\textbf{more sensitive to this area's data than to any other area in Hull}
(rank 1 of 168): if this area's information were replaced with typical Hull
values, its predictions would change more than anywhere else. Neighbouring
areas share this pattern: together with the rest of the Old Town and the
areas running west along Hessle Road and Anlaby Road towards Dairycoates
(around Coltman Street), north-west along Spring Bank into Stepney (around
Argyle Street), and north along Beverley Road and Freetown Way, it forms a
group of nine areas where the model is unusually sensitive to local data.

This describes what the model uses, \textbf{not what causes crime}, and says
nothing about the people who live, work or shop there.
